\section{Vulnerability Replay Evidence and Rediscovery Protocol}
\label{sec:cve}
The vulnerability study comprises four cases: three associated with CVE identifiers and one SMARTPL availability failure without a CVE. Each case has a report and a triggering input; three have direct execution evidence, and ProFTPD also has a reported patched-version comparison. \namedref{Tables}{tab:cve} and~\ref{tab:cveresults} summarize these observations. Case identities, versions, and replay counts follow the original reports; incomplete patch and experimental-configuration evidence prevents independent verification of every version pair.

\subsection{Case identity and evidence provenance}
We distinguish a direct execution observation (L) from a replay result stated only in a case report (R). A report can contain additional trials not represented by the available execution evidence. Duplicate diagnostic observations do not constitute independent trials. Input identity and byte length are retained in the supporting evidence. This retrospective analysis involves no additional target executions. \namedref{Table}{tab:cve} identifies the reported versions, triggering conditions, and comparative evidence.

\begin{table*}[!tbp]
\centering\small
\caption{Vulnerability cases and supporting evidence (L: direct observation; R: report).}
\label{tab:cve}
\begin{tabularx}{\textwidth}{@{}>{\raggedright\arraybackslash}p{0.20\textwidth} >{\raggedright\arraybackslash}p{0.19\textwidth} Y Y@{}}
\toprule
Case / interface & Reported vulnerable version & Trigger evidence / oracle & Patch or control evidence \\
\midrule
\mbox{CVE-2023-51713}\newline ProFTPD / FTP & Reported vulnerable build; 65,535-byte buffer. & 60,006-byte input; ASan heap-buffer-overflow (L). & Security patch: three rejections, no ASan (L). \\
\addlinespace[3pt]
\mbox{CVE-2026-38998}\newline LIVE555 / RTSP & Campaign: 2023.05.10; report: 2026.02.26 (R). & 728-byte public PoC; use-after-free (L). & Local patch: 1/1 triggers (R); upstream comparison unavailable. \\
\addlinespace[3pt]
\mbox{CVE-2026-39863}\newline Kamailio / SIP over TCP & Reported vulnerable build; TCP required. & 278-byte input; parsed Content-Length $-10$ (L). & Patched: 0/253 negative-value outcomes; guard observed (L). \\
\addlinespace[3pt]
SMARTPL availability case\newline forked-daapd / HTTP & Reported version 27.2. & 148-byte input; unresponsive service, live process (R). & No patched replay/CVE; removal reported in 28.4. \\
\bottomrule
\end{tabularx}
\end{table*}

For ProFTPD, the reported defect is an out-of-bounds read during FTP command processing, supported by an ASan heap-buffer-overflow observation. The 60,006-byte reproducer caused a one-byte read beyond a 65,568-byte region. The oracle uses a 65,535-byte command buffer; the report states that the default 512-byte configuration does not expose the same error to ASan. The result is therefore conditional on this setting. The full replay in \namedref{Table}{tab:cveresults} uses a purpose-built version pair at the campaign's base commit whose only difference is the minimal upstream fix; on that pair, 172 of 233 campaign crash inputs trigger the \texttt{make\_ftp\_cmd} predicate on the vulnerable side and none of 242 activate ASan on the patched side. This supplies differential replay evidence with verified version identity for the pair, while the earlier three-trial fix batch remains separate. Earlier reported replays yielded 1/1 seed, 1/1 fuzzer-crash, and 3/3 batch-crash reproductions; the separate fix batch yielded 0/3 using the same input and buffer.

For LIVE555, direct diagnostic evidence supports a use-after-free triggered by the public PoC. A duplicated observation is treated as one diagnostic signature rather than independent replication. The report attributes fuzzer-generated inputs to the same memory-lifetime failure. It gives 1/1 public-PoC reproduction, 3/3 repeats for each of three fuzzer inputs, a later 8/9 aggregate, and 1/1 persistence after a local patch. The fuzzer-generated reproducer, complete memory-lifetime evidence, and individual trial outcomes are unavailable. The public PoC supports replay validation, but does not demonstrate independent rediscovery by the evaluated controller or independently establish cross-version CVE attribution.

For Kamailio, a negative parsed Content-Length supports the reported integer-overflow condition. The trigger requires TCP and is not assessed by the UDP-only benchmark condition. The full census in \namedref{Table}{tab:cveresults} replays every queue entry carrying an oversized Content-Length on a purpose-built pair at the campaign's base commit whose only difference is the upstream guard: 102 of 253 entries produce the negative-value artifact on the vulnerable side, none do on the patched side, and the patched build logs the upstream guard for 193 of 253. Earlier reported batches remain separate: 1/1 seed and 3/4, 19/20, 19/20, and 50/51 selected inputs. The prior local-guard observation is also separate. Queue entries are fuzzer-generated artifacts, so the census measures predicate yield over retained inputs, not rediscovery. An earlier unsupported 20/20 versus 0/20 claim is excluded from quantitative results. 

For SMARTPL, the 148-byte input reportedly causes lexer errors and a failed health probe while the process remains alive. Reported repeats comprise 5/5 advisory trials, 3/3 trials with a 48-message seed, and a later 1/1 replay. None of the five recorded campaign hangs reproduces this condition; a CVE and verified patched counterpart remain unavailable. Parser removal in version 28.4 is reported, not a verified patched replay.

\subsection{Replay outcomes and their denominators}
\namedref{Table}{tab:cveresults} summarizes the replay outcomes; earlier reported batches remain separate in the case descriptions. A numerator counts replay attempts or selected queue inputs satisfying that case's predicate; the denominator is the number tested. These are not independent 12-hour fuzzer runs, and the main experiment's nominal $N=10$ is not substituted for the actual replay denominator. Selected batches, per-input repeats, and later aggregate counts may overlap; they are not pooled into one success rate. In particular, 19/20 selected inputs is not 95\% campaign recall, and 8/9 LIVE555 replays is below the prospective 95\% stability threshold.

Three complete-input replays extend these batches (September 30 observations). First, every crash input recorded by the two 24-hour ProFTPD LoopFuzz campaigns (233 in-campaign inputs: 121 and 112 per run; 9 teardown inputs) was replayed once against the vulnerable and patched images of the verified \texttt{o-pftp-01} pair (same base commit, minimal upstream fix). On the vulnerable image, 172 of the 233 in-campaign inputs trigger the ASan heap-buffer-overflow predicate in \texttt{make\_ftp\_cmd}; all 172 carry that frame and no other signature appears. The patched image shows no ASan activation for any of the 242 inputs. The 61 non-reproducing inputs and the 9 teardown inputs may require campaign state that a standalone replay does not reconstruct; the in-campaign forensic sidecars recorded almost no usable ASan output, so the pair replay is the authoritative verdict. Second, all 88 LIVE555 campaign crash inputs were replayed against the campaign build, which contains no sanitizer: 82 terminate the server and 6 do not. This establishes reproduction on that build only; without a patched counterpart or sanitizer output it does not attribute the crashes to the use-after-free. Third, all 253 Kamailio queue entries with Content-Length $\geq 2^{31}$ were replayed on the \texttt{o-kam-01} pair: 102 produce the negative-value artifact on the vulnerable build and none on the patched build. The two runs contributed 77 and 176 entries, with yields of 36/77 and 66/176. Overflow values include $-2147483648$, $-1$, and $-954437177$; 151/253 entries do not produce the artifact, with ordering and parse-path explanations unresolved. All three replays are artifact-level measurements: they provide no arm-level recall, time-to-trigger, or cost denominators.

\begin{table*}[!tbp]
\centering\small
\caption{Replay outcomes and controls; original denominators retained.}
\label{tab:cveresults}
\begin{tabularx}{\textwidth}{@{}>{\raggedright\arraybackslash}p{0.20\textwidth} Y Y Y@{}}
\toprule
Case & Replay outcome & Control outcome & Evidence limit \\
\midrule
\mbox{CVE-2023-51713} & 172/233 ASan triggers (L). & Patched: 0/242 ASan activations (L). & Artifact replay; campaign recall unmeasured. \\
\addlinespace[3pt]
\mbox{CVE-2026-38998} & 82/88 server terminations (L). & Local patch: 1/1 triggers (R). & No sanitizer or verified patched build. \\
\addlinespace[3pt]
\mbox{CVE-2026-39863} & 102/253 negative-value outcomes (L). & Patched: 0/253; guard: 193/253 (L). & Selected queue entries; campaign recall unmeasured. \\
\addlinespace[3pt]
SMARTPL availability case & Advisory: 5/5; seed: 3/3; later: 1/1 (R). & Campaign hangs: 0/5 matching outcomes (R). & Known-input replay; no patch or CVE. \\
\bottomrule
\end{tabularx}
\end{table*}

Additional reported crashes and campaigns cannot be linked to the primary benchmark and ablation observations. Crash totals do not count unique vulnerabilities, and durations such as 60 or 1,540 minutes do not specify time to first trigger. Associations among reproducing inputs, verified A--E conditions, model information exposure, and independent root causes are incomplete. The replay outcomes therefore do not establish arm-specific discovery probabilities, novel CVEs, or admission/calibration gains.

\subsection{Controlled rediscovery still to be measured}
Confirmatory rediscovery requires a vulnerable version and the same version with a minimal security patch, holding dependencies, instrumentation, settings, and input compatibility constant. Each case needs a network-reachable trigger, independent reached/triggered criteria, reproducible initial conditions, and a validated patch comparison. The proposed six-case pilot across at least three protocols, followed by expansion toward twelve, remains incomplete. Eligibility requires at least 95\% observed oracle repeatability in a declared validation set, no activation after patching, reproducible experimental preparation, and information separation. The four cases do not constitute six qualified version pairs; a small all-success replay batch does not establish an underlying success probability of at least 95\%.

Information-controlled campaigns must exclude CVE identifiers, advisories, public PoCs, and patch contents from model inputs and campaign seeds, while retaining anonymous oracle IDs in separate benchmark metadata. The public PoC and known triggers in this collection are validation inputs; their replay is not evidence of information-controlled rediscovery. Source-level trigger conditions, sanitizer reports with patch validation, and protocol invariants distinguish reaching relevant code from actually triggering a defect. Experiments require isolated target networks and reset between trials; these are protocol requirements, not properties inferred from the reports.

For each CVE and arm, the formal plan retains ten independent runs and a 12-hour deadline, prespecified before execution. Deadline recall, time-to-trigger, reached-but-not-triggered counts, patched-target activation, executions/coverage/tokens before trigger, and recall per CPU-hour or 100,000 tokens remain unavailable for the controlled A--E comparison. Nontriggers require right-censoring records, and interrupted runs require separate failure accounting. The available replay batches do not define those risk sets, so Kaplan--Meier estimates or arm-level security-effectiveness claims are not reported.
